% The theme switch, and everything shared across themes.
%
% ---- SWITCH THE DECK'S THEME BY EDITING THE ONE LINE BELOW ----------------
% themes/ holds one file per theme: the palette, the symbolic bindings the
% slides are written against, and \themelogo. Nothing else in the deck names a
% palette colour, so this line is the whole change.
% anthropic_orange — palette and symbolic bindings. \input by colors.tex.
%
% After ~/DavidFillmore/Media/templates/anthropic_orange/colors.tex, with ONE
% deliberate departure. That template is single-seam: it sets
% accentSoft = accent, so every tier is the one clay. This deck cannot use
% that. Slides 02 and 04 are built on a real contrast between a block and an
% alertblock — specification against constraint, archive against trap — and
% under a single seam the two blocks become the same colour and the argument
% goes with them. So accentSoft is bound to Anthropic's dark brown, the other
% end of the same brand family: the alert tier reads as the grave one rather
% than the loud one, which suits a warm palette with no red in it.

% --- Anthropic palette ---
% Clay is extracted from Claude_by_Anthropic.svg; the browns are the family it
% sits in, as recorded in Media/decks/agentic_ai/colors.tex.
\definecolor{claudeClay}{RGB}{217,119,87}    % Anthropic orange
\definecolor{claudeBrown}{HTML}{3D3929}      % dark brown
\definecolor{claudeKraft}{HTML}{D4A27F}      % kraft, the light end
\definecolor{claudeBrownMid}{HTML}{886D54}   % midpoint of brown and kraft

% --- Symbolic seam ---
% Slides use ONLY these; never the claude* names above.
%
% CONTRAST, measured, not assumed: white on clay is 3.12:1. That clears WCAG AA
% for large text and does not clear it for body text, so clay is a GROUND for
% bold headings only — the frametitle bar and block titles — and never carries
% body copy. Body-weight text takes accentDeep (11.9:1 on white) or accentSoft
% (11.0:1). This is the same trade the shipped anthropic_orange template and
% agentic_ai's \claudeslide both make, and it is fine on a projector at size;
% it is written down here so it stays a decision rather than an oversight.
\colorlet{accent}{claudeClay}                 % primary chrome
\colorlet{accentSoft}{claudeBrown}            % second tier: alerts, annotation
\colorlet{accentDeep}{claudeClay!62!black}    % deep ink: diagram text, captions
\colorlet{accentFill}{claudeKraft!32!white}   % diagram node fill
\colorlet{accentTint}{claudeClay!10!white}    % block body ground
\colorlet{mutedInk}{claudeBrown!52!white}     % de-emphasised labels

\newcommand{\themelogo}{Claude_symbol_clay.pdf}

% % nasa_blue — palette and symbolic bindings. \input by colors.tex.
% After ~/DavidFillmore/Media/templates/nasa_blue/colors.tex, extended with the
% three seam names this deck's slides need beyond the shipped accent trio.

% --- Official NASA brand palette ---
% NASA's core brand is three colours (blue, red, white), stable since 1992.
% NASA Blue PMS 287 C, NASA Red PMS Bright Red C. White is the third — use
% plain white.
\definecolor{nasaBlue}{HTML}{0B3D91}
\definecolor{nasaRed}{HTML}{FC3D21}

% --- Symbolic seam ---
% Slides use ONLY these; never the nasa* names above. That is the portability
% contract that makes a theme switch a one-line edit in colors.tex.
\colorlet{accent}{nasaBlue}              % primary chrome: title bar, block titles
\colorlet{accentSoft}{nasaRed}           % second tier: alerts, annotation, gates
\colorlet{accentDeep}{nasaBlue!70!black} % deep ink: diagram text, captions
\colorlet{accentFill}{nasaBlue!16!white} % diagram node fill
\colorlet{accentTint}{nasaBlue!8!white}  % block body ground
\colorlet{mutedInk}{black!42}            % de-emphasised labels

% The meatball is a full-colour insignia and does not reverse to white by a
% fill swap; on a dark ground it needs a white plate.
\newcommand{\themelogo}{nasa_logo.pdf}

% --------------------------------------------------------------------------
%
% The seam the slides may use, and the only one:
%   accent      primary chrome — frametitle bar, block titles. Bold text on it
%               only; see the contrast note in themes/anthropic_orange.tex.
%   accentSoft  second tier — alertblocks, annotation, review gates
%   accentDeep  deep ink — diagram labels, \figcaption, body-weight emphasis
%   accentFill  diagram node fill
%   accentTint  block body ground
%   mutedInk    de-emphasised labels (artifact counts under the phase chain)
%   \themelogo  the theme's own mark
%
% A slide that reaches past these for a literal colour name has broken the
% contract and will not survive the next switch. Check with:
%   grep -ohE 'nasa[A-Za-z]+|claude[A-Za-z]+' frames/*.tex

% --- Vendor inks, deliberately NOT part of the seam -----------------------
% The CLI transcript panel paints Claude Code's own interface: the clay `>`
% prompt and the clay bullet are that product's colours, the way the NASA
% meatball is NASA's. They are content, not chrome, and must not follow the
% theme — under nasa_blue the panel still has to look like the terminal it is.
% Defined here rather than in a theme file for exactly that reason.
% (anthropic_orange defines claudeClay too, to the same value; a second
% \definecolor of a name already defined is a redefinition, not an error.)
\definecolor{claudeClay}{RGB}{217,119,87}

% --- Font binding ---------------------------------------------------------
% Helvetica (helvet is the shipped stand-in), kept across both themes even
% though Media's anthropic_orange template specifies Latin Modern Sans. Two
% reasons: LM Sans sets about 15% narrower, so every hand-tuned column width,
% figure height and \vspace in frames/ would need re-measuring; and the
% plots carry their own Helvetica-like sans, which a Latin Modern deck would
% sit oddly against. \familydefault is set explicitly because nothing else
% picks sans by default.
\usepackage{helvet}
\renewcommand{\familydefault}{\sfdefault}

% --- Graphics paths (relative to slides/) ---------------------------------
% Both asset directories live INSIDE the deck. The marks were copied out of
% ~/DavidFillmore/Media/graphics/logos/ rather than referenced across the
% filesystem, so this repo builds the deck on its own: a relative path out to
% another checkout would break for anyone who cloned only earth-eval.
% Provenance and usage rules: graphics/logos/SOURCES.md.
\graphicspath{%
  {figures/}%
  {graphics/logos/}%
}
